\section{Experimental Setup}
\label{sec:setup}

We evaluate HESOD on two demanding high-resolution remote sensing benchmarks:
\begin{itemize}
    \item \textbf{UAVDT} \cite{UAVDT}: An urban aerial vehicle benchmark evaluated on the official test split across three classes (\emph{car}, \emph{truck}, and \emph{bus}) under an input resolution of $1280 \times 1280$ pixels.
    \item \textbf{SeaPerson (TinyPersonV2)} \cite{TinyPerson,SeaDronesSee}: A maritime search-and-rescue benchmark comprising 300,375 annotated test instances across 5,752 images, with over 85\% of targets falling into the \emph{Very Tiny} and \emph{Tiny} categories ($<32\times 32$ pixels). We process full whole-frame images at $2048 \times 2048$ resolution without artificial pre-cropping or sliding-window tiling.
\end{itemize}

\textbf{Implementation and Training.} All models are implemented in PyTorch using a standard YOLOv5m architecture as the shared base detector for fair comparison. Models are trained for 50 epochs using SGD with cosine annealing, following standard selective detection training protocols \cite{ESOD}.

\textbf{Evaluation Metrics.} Detection accuracy is measured using standard COCO-style metrics ($\mathrm{mAP@.5}$ and $\mathrm{mAP@.5:.95}$). To rigorously evaluate scale sensitivity, we report size-stratified detector recall across four physical pixel bins: \emph{Very Tiny} ($<16^2$ px), \emph{Tiny} ($16^2-32^2$ px), \emph{Small} ($32^2-96^2$ px), and \emph{Medium/Large} ($>96^2$ px). Efficiency is evaluated via total GFLOPs, frame rate (FPS), Bounding Patch Recall (BPR), and spatial retention ratio ($\mathrm{Occupy}$).
